% abstract.tex

\begin{abstract}
No-reference mean opinion score (MOS) predictors output a single number with no signal of trustworthiness, a risk in automated pipelines. We test whether distance from a reference feature distribution can supply that signal, using three frozen predictors, DNSMOS P.835, NISQA, and UTMOS, in a leakage-controlled design that separates feature-reference fitting, content-disjoint calibration, and evaluation across all seven NISQA Corpus subsets. Euclidean, Mahalanobis, and nearest-neighbor distances detect which corpus a clip came from (mean domain AUROC 0.615--0.690) but do not flag which clips carry the largest prediction error within a corpus (mean AUROC 0.433--0.499, near chance). Split conformal intervals built on these distances reach 90.1--91.8\% coverage at the 90\% nominal level with widths of 2.26--2.91 MOS on the four-point scale. Stratifying interval width by Mahalanobis distance helps some predictor--domain pairs and hurts others. This pattern reproduces on an independent multilingual speech-enhancement corpus, so practitioners should read a feature-distance alarm as a domain-mismatch flag, not a per-utterance reliability score.
\end{abstract}
